% !PRINTED 227
\weberpage{227}
Der Ausdruck auf der rechten Seite von (7) ist eine \gesperrt{bilineare}
Function der $t$ und $\tau$; er möge entwickelt die Form
haben
\[
\sum_{0,\,n-2}^{h} \sum_{0,\,n-2}^{k} B_{h,k} t_h \tau_k,
\]
mit der Bedingung $B_{i,k} = B_{k,i}$. Die \gesperrt{Bezoutiante} ergiebt
sich dann, wenn man $\tau = t$ setzt, also
\weq{8}{B = \sum_{0,\,n-2}^{h,k} B_{h,k} t_h t_k,}
worin die Coëfficienten $B_{i,k}$ quadratische Functionen der $a$ sind;
und aus (7) folgt
\weq{9}{S(y z) - \frac{1}{n} S(y) S(z) = \sum_{0,\,n-2}^{h,k} B_{h,k} t_h \tau_k,}
\weq{10}{S(y^2) = \frac{1}{n} [S(y)]^2 + B.}
Nach der Formel (8), \S.\,68 ist
\[
\begin{aligned}
y - \frac{1}{n} S(y) &= t_{n-2} F_0 + t_{n-3} F_1 + \cdots + t_0 F_{n-2} \\
z - \frac{1}{n} S(z) &= \tau_{n-2} F_0 + \tau_{n-3} F_1 + \cdots + \tau_0 F_{n-2},
\end{aligned}
\]
und wenn wir beides multipliciren,
\[
y z - \frac{1}{n} [y S(z) + z S(y)] + \frac{1}{n^2} S(y) S(z)
= \sum_{0,\,n-2}^{h,k} \tau_h t_k F_{n-h-2} F_{n-k-2}.
\]
Summiren wir diese Formel über alle Wurzeln $x$, d.\,h. nehmen
wir die mit $S$ bezeichnete Summe, so folgt
\[
S(y z) - \frac{1}{n} S(y) S(z) = \sum^{h,k} \tau_h t_k S(F_{n-h-2} F_{n-k-2}),
\]
und die Vergleichung mit (9) lehrt
\weq{11}{B_{h,k} = S(F_{n-h-2} F_{n-k-2})}
Wir bezeichnen mit $\varDelta$ die Determinante der Function $B$,
setzen also
\[
\varDelta = \sum \pm B_{0,0} B_{1,1} \ldots B_{n-2,\,n-2}.
\]
Nach dem Multiplicationssatz der Determinanten und mit
Rücksicht auf die Formeln
\[
S F_0 = 0, \quad S F_1 = 0, \ldots S F_{n-2} = 0
\]
%CONT>>
